\chapter{方法论收束：证据分级、可复现性与研究伦理}\label{ch:12}

\begin{nlkeypoints}
  \begin{enumerate}[label=(\arabic*),leftmargin=2.4em]
    \item 全书的证据分级制度（\conf{A}/\conf{B}/\conf{C}/\conf{M}）
          在本章给出设计理由、操作细则与三条局限；
          它是本书的\keyw{约定}而非行业标准，
          其价值在执行的一致性，不在分级的权威性。
    \item 本章交付“\keyw{可复现性清单}”——
          把散落各章的纪律（口径并列、E1--E3、
          措辞模板、模型声明、编码公开）收拢为
          修订者可直接执行的十项检查
          （第三、四、六章承诺的兑现）。
    \item 第一章遗留的元问题在此收束：
          当研究对象以“解构严肃性”为乐时，
          严肃教材的自洽出路不是假装对象严肃，
          而是\keyw{让体例承担幽默、让方法承担严肃}——
          语域错配因此不是噱头，而是方法立场
          （\ref{sec:02:legitimacy} 的最终答案）。
  \end{enumerate}
\end{nlkeypoints}

全书以个案始（第一章的现象描述），
以方法终。这样安排不是惯例使然，而是因为
奶蛙案例真正可迁移的部分，恰恰不是“奶蛙是什么”
（它是速朽的），而是\keyw{研究一个正在变异、
正在消失、正在被商业与权威觊觎的网络对象时，
研究者应当遵守什么规程}。本章把散落全书的
方法装置收拢为三组：证据分级（第一节）、
可复现性清单（第二节）、伦理三条线（第四节），
并在第三节处理网络材料的时效问题，
在第五节回答那个“教材自身立场”的元问题。

\section{证据分级制度：设计、执行与局限}\label{sec:12:grades}

\subsection{设计理由}

\conf{A}/\conf{B}/\conf{C}/\conf{M} 四级
（第一章 \ref{sec:01:method}）的区分维度是
\keyw{来源的产生方式}，不是内容的真伪：
A 级是职业编辑流程与署名担保的产物，
B 级是机构信誉担保的公开信息，
C 级是社群生产、无编辑流程的材料，
M 级是本书为教学构造的人工物。
用产生方式而非真伪作区分维度，理由有二：
其一，真伪通常无法当场判定，
产生方式却可以当场核验；
其二，\keyw{真伪判断必须由读者带着分级信息自行完成}，
制度不应替读者下结论——这是本书与
“本人已核实一切”式写作的分野。

\subsection{操作细则的五个要点}

第一，\keyw{级别跟随陈述而非来源}：同一来源的
不同陈述可以标不同级别（第三章“上线日期”取
\conf{B} 级倒推即是实例）。第二，\keyw{多口径并列}：
同一指标的不同来源数值并列呈现，不合并、不取均值
（表 \ref{tab:03:numbers}）。第三，\keyw{C 级不得作为
论证前提}：C 级材料可以引出问题、提供线索、
充当被检验对象，但结论不得悬于其上——
第七章 H7.5 的证伪条件设计即遵守此条。第四，
\keyw{M 级必须显式自首}：模型、示例数据、教学命题
在首次出现处声明构造身份，防止教学装置
在二次引用中洗白为经验发现。第五，
\keyw{措辞随级别走}：“据……报道”（A/B）与
“据社群考据”“有说法认为”（C）的句式差异
是分级制度在语言层的执行（各章实例从略）。

\subsection{三条局限}

\begin{enumerate}[label=\arabic*.]
  \item \textbf{级别是粗粒度的}：同为 \conf{C} 级，
        失传媒体维基的词条（有引用、有修订史）与
        无出处的水贴证据地位不同，四级制无法表达
        这一梯度。使用中的补救是\keyw{在正文里
        保留材料自身的限定语}（第四章“不替考据者
        删掉他们自己写的‘但是’”）；
  \item \textbf{分级判断本身可错}：把一则伪装成
        官方通稿的营销稿标为 \conf{B} 级，
        制度无能为力。分级制度降低错误的传导率，
        不保证零错误；
  \item \textbf{制度有成本}：严格的 C 级纪律
        意味着大量“有趣但不可用”的材料被搁置
        ——本书对奶蛙周边的具体价格、销量
        一字不提，就是这一纪律的代价。
        我们认为这个代价值得付，
        但修订者应知道它是\keyw{选择}的结果。
\end{enumerate}

\section{可复现性清单}\label{sec:12:checklist}

以下是散落各章纪律的收拢版，供修订者
在交付新版本前逐项自查（第三章、第四章、第六章
承诺的“操作化版本”）：

\begin{table}[htbp]\centering
  \caption{可复现性清单（本书构造，\conf{M} 级）}
  \label{tab:12:checklist}
  \small
  \begin{tabularx}{\textwidth}{@{}clXl@{}}
    \toprule
    序号 & 检查项 & 操作要点 & 来源 \\ \midrule
    1 & 引用可追溯 & 每条事实有 \conf{A/B} 级出处或明确标注 \conf{C} &
      第一章 \\
    2 & 口径并列 & 同一指标多来源数值并列，不合并 &
      第三章 \\
    3 & C 级措辞 & “据社群考据 / 有说法认为”句式 &
      第一章 \\
    4 & 自我限定保留 & 不删改考据材料自带的“但是” &
      第四章 \\
    5 & E1 对象限定 & 考据叙述宾语限于账号、帖子等公开对象 &
      第四章 \\
    6 & E2 推测隔离 & 个人境况推测“说明存在、不转述内容” &
      第四章 \\
    7 & E3 权利保留 & 二创形象叙述附署名权归属 &
      第四章 \\
    8 & M 级自首 & 模型 / 示例数据首次出现处声明构造身份 &
      第五章 \\
    9 & 方法可复现 & 判定标准、样本框、计数规则写定 &
      第八章 \\
    10 & 编码与量表公开 & 编码手册、问卷、权重随文交付 &
      第十一章 \\ \bottomrule
  \end{tabularx}
\end{table}

清单的用法说明两点。其一，清单是\keyw{拦截装置}：
第二章警告过，语域错配的严肃性不能替代证据的严肃性，
“若某处用 C 级考据支撑了 A 级口吻的结论，
错配就从立场表达滑向误导”——第 3、4 项正是
为拦截这种滑动而设。其二，清单是\keyw{最低标准}：
通过十项不等于研究合格，违反任何一项则肯定不合格。
修订者可以在其上追加本校条款（如增设校内案例的
核对规程），但不得删减既有条目——删减须在新版
前言中显式声明并说明理由。

\section{网络材料的时效：引文腐烂与现存样本}\label{sec:12:rot}

第四章把“不可复核性”处理为个案
（智能体 \fname{zNCvQDKfXy} 的消失），
本节把它一般化。网络引文面临系统性的
\keyw{引文腐烂}（\en{link rot}）：链接失效、
内容被编辑、账号注销、平台下架。
本书的三项应对及适用边界：

\begin{enumerate}[label=\arabic*.]
  \item \textbf{访问日期制度}：所有网络来源记录
        首次访问日期（附录 D），读者以该日期为
        内容状态的基线。边界：访问日期只证明
        “当时如此”，不担保“现在如此”；
  \item \textbf{现存样本立场}：一切基于可检索内容的
        分析，对象都是\keyw{幸存下来的样本}——
        已删除、已注销、未备份的部分永久缺席
        （第十章“样本一端在消失、一端在缺席”）。
        边界：该立场要求结论措辞留有
        “在现存材料中”的限定，本书执行不一贯处，
        以本条为准修订；
  \item \textbf{不可复核性的显式登记}：凡关键证据
        载体已消失的指向（如奶蛙起源），必须在
        结论处写明“不可复核”，不得以交叉印证
        冒充复核。边界：登记不改变指向的
        证明力，只防止证明力被高估。
\end{enumerate}

宏观的网民与短视频使用口径应引用权威统计的
当期报告并注明报告期数\cite{baidu-index}，
避免以过期口径描述现状——这是 \conf{B} 级材料
在时效维度上的延伸纪律。

\section{研究伦理的三条线}\label{sec:12:ethics}

\subsection{对象伦理：E1--E3}

第四章已给出 E1（对象限定）、E2（推测隔离）、
E3（权利保留）的完整表述，此处只补一点
为什么它们必须升格为全书规则而非第四章的
地方性纪律：奶蛙案例中，考据完成与舆论要求
“当事人出面”之间只有一步，任何转述个人境况的
文字——哪怕以同情的形式——都在给这一步施压。
第十一章把 E1--E3 接入访谈设计（以角色为宾语），
第十章的赞诗分析对其只作语体分析、不复制全文，
都是同一规则在不同材料上的执行。

\subsection{版权伦理：分析与背书的界线}

本书把未经授权的形象再生产作为\keyw{分析对象}
（第三、八、九章），附录 D 申明了使用限于
学术评论与教学示例。方法层面补一条操作界线：
\keyw{分析陈述的宾语是结构，不是行为人}——“
换皮生产降低了变异成本”是分析，“你应该去换皮”
是背书，两者的句式区别与 E1 同构。
凡是发现本书某处从结构滑向了怂恿，
按可复现性清单第 1 项的标准，那是缺陷。

\subsection{社群伦理：研究者与社群的互惠}

研究依赖社群产生的材料（考据词条、第一版教材、
同好会记录），却常以“中立观察者”自居。
本书的立场：\keyw{依赖即负债}。可行的偿还方式包括：
引用时忠实署名（C 级材料的来源标注即对
考据者劳动的承认）、研究结果向社群回馈
（第十一章伦理条款第 5 条）、以及不利用
信息差损伤社群（例如不公开尚未被社群公开的
考据线索）。孙立平式的宏观社会心态分析
\cite{sunliping2010} 为“自嘲式表达”提供的
背景框架，也应向受访者开放讨论——
受访者有权知道研究者用何种理论框定他们的笑。

\section{语域错配的最终答案}\label{sec:12:register}

第一章遗留了全书最深的元问题：
当研究对象本身以“解构严肃性”为乐时，
一本严肃教材如何保持自身立场的一致性？
第二章的答案（\ref{sec:02:legitimacy}）是：
命名与体例的语域错配是自觉的方法表达，
但其严肃性不能替代证据的严肃性。
本节给出最终版本，分三点。

第一，\keyw{幽默的安放位置是体例，不是论证}。
章名、框标题、脚注可以玩笑，因为它们不承担
证明功能；定义、数据、推理链必须中性，因为
证明功能不可戏仿。这一分工在本书的执行是：
所有彩色框承担“人声”，所有编号环境
（定义、命题、假设）承担“制度声”。
第二，\keyw{对象的无厘头性是数据，不是许可}。
奶蛙的“意义不明”是第六章的分析对象，
不是本书放松论证标准的理由——恰恰相反，
对象越混沌，方法越要清晰，否则分析
不过是给混沌再添一层。第三，
\keyw{严肃的最终检验是可批判性}：
本书自始至终把可被读者纠正当作目标状态
（第一章：“读者若发现某条 C 级陈述被当作
事实使用，那正是本书鼓励你指出的错误”）。
一本鼓励读者推翻自己的教材，
它的严肃不需要靠板着脸维持。

\begin{nlcritique}[批判性反思：对本书自身的最后检验]
  以自己的标准检验自己：全书是否处处通过了
  表 \ref{tab:12:checklist}？诚实回答是
  \keyw{没有完全通过}。已知的薄弱处：
  其一，第二章“抽象文化”定义标注了 \conf{M} 级，
  但后续章节调用它时偶有接近当作既定概念的表述；
  其二，第十章的协商强度谱三维读数依赖单一来源，
  读数本身应视为“待复核的示例”；
  其三，各章对 \conf{C} 级材料的措辞纪律
  执行不一——个别句子仍用陈述句式引出了
  本应“据社群考据”的内容。这些薄弱处
  不掩盖、不辩解，留作修订第一优先级。
  一本教人识别证据的方法书，
  最不能做的事就是暗示自己没有病灶。
\end{nlcritique}

\section{面向修订者}\label{sec:12:revisers}

本书以 CC BY-NC-SA 4.0 发布（附录 D），
修订与续编是被预期的使用方式。操作规程：

\begin{enumerate}[label=\arabic*.]
  \item \textbf{事实手册先行}：新版修订从更新
        事实手册开始——核对每条 \conf{A/B} 级口径的
        现时状态、重新访问 \conf{C} 级来源并更新日期、
        把新事件补入年表（附录 A）；
  \item \textbf{分级制度保留}：新版可以修正分级、
        补充级别，但不得废除四级制或删除
        表 \ref{tab:12:checklist} 的条目
        （\ref{sec:12:checklist}）；
  \item \textbf{模型重做自由}：SIR 模型、协商强度谱、
        语音距离指标均可替换参数、权重甚至整体架构，
        保留的约束只有一条——“每个构件都声明
        自己的教学身份”；
  \item \textbf{本校化}：补充本校案例时遵守
        第十章的 C 级纪律与 \ref{sec:12:ethics}
        三条线；校本材料（贺图、群聊记录）
        须获当事人同意后方可收录；
  \item \textbf{版本记录}：每次修订在前言后附
        修订说明（改了什么、为什么、依据哪条清单项），
        使读者能追踪本书自身的演变——
        一本关于“考据与版本”的书，
        应当对自己的版本负责。
\end{enumerate}

\section{本章小结}\label{sec:12:summary}

本章把全书的方法装置收拢为三组并交代其边界：
证据分级制度（产生方式为区分维度、五条操作细则、
三条诚实的局限）、可复现性清单（十项拦截装置，
最低标准）、时效纪律（访问日期、现存样本立场、
不可复核性登记）；伦理部分把 E1--E3、
版权界线、社群互惠并为三条线；
最后以三点收束了第一章的元问题——
幽默安放于体例、无厘头性是数据而非许可、
严肃的检验是可批判性——并以自我检验的方式
列出了本书自身的三处病灶。全书在此结束：
十二章走过了一个现象的描述（一、三章）、
考据（四章）、工具（五章）、解剖（六章）、
机制（七章）、生产（八章）、经济（九章）、
空间（十章）、设计的交付（十一章）与方法的自审
（本章）。它没有解决奶蛙为什么好笑——
那要留给读者自己的笑；它解决的是
\keyw{谈论这种好笑时，怎样谈论才算数}。

\begin{nlexercises}
  \item 从本书正文挑出一个你认为\keyw{违反了}
        可复现性清单某一项的具体句子
        （第十二章批判框列出的三处病灶是提示，不是上限），
        写出：清单条目号、原句、修改后的句子、
        修改依据。
  \item 设想你所在的高校同好会决定续编本书并补充
        本校立牌事件。按 \ref{sec:12:revisers} 的规程
        写出续编的前三步操作（含每步的伦理检查点），
        并指出哪一步最容易省略、省略的后果是什么。
  \item “幽默的安放位置是体例，不是论证”
        ——为这句话找一个\keyw{反例}：
        一部（任意领域的）严肃作品，
        其幽默进入了论证本身且因此更有说服力。
        分析该例成立需要什么前提条件，
        并说明这些条件在奶蛙研究中是否具备。
\end{nlexercises}

\nlkeywords 证据分级、可复现性清单、引文腐烂、
现存样本、对象伦理、版权界线、社群互惠、
语域错配、修订规程、自我检验
